\documentclass{article}
\usepackage{listings}
\lstdefinelanguage{JavaScript}{
  keywords={break, case, catch, continue, debugger, default, delete, do, else, finally, for, function, if, in, instanceof, new, return, switch, this, throw, try, typeof, var, void, while, with},
  basicstyle = \ttfamily,
  sensitive=true,
  comment=[l]{//},
  morecomment=[s]{/*}{*/},
  morestring=[b]", 
  morestring=[b]'
}
\usepackage{color}
\usepackage{graphicx}

\title{RefSynの手動解析記録}
\author{}
\date{}

\begin{document}

\maketitle

\section{対象関数: \texttt{append}}
現在のノードを先頭とする連結リストの末尾に、新しいノードを追加するメソッド。
引数として渡された値を \texttt{val} フィールドにもつノードを新たに作成し、それをリストの最後に接続。

\subsection{初期プログラム}

\begin{lstlisting}[language=JavaScript]
class Node {
  append(arg) { /* TBD */ }
}
var lst = new Node(); 
lst.val = 2;
lst.append(0);
lst.append(3);
\end{lstlisting}

\subsection{操作ログ}

\subsubsection{\texttt{lst.append(0)} に対する操作}
\begin{itemize}
  \item \texttt{addNode}: \_\_temp1, Node, false, undefined
  \item \texttt{addNode}: \_\_temp2, 0, true, string
  \item \texttt{addEdge}: \_\_temp1, \_\_temp2, val
  \item \texttt{addEdge}: main-new1, \_\_temp1, next
\end{itemize}

\subsubsection{\texttt{lst.append(3)} に対する操作}
\begin{itemize}
  \item \texttt{addNode}: \_\_temp3, Node, false, undefined
  \item \texttt{addNode}: \_\_temp4, 3, true, string
  \item \texttt{addEdge}: \_\_temp3, \_\_temp4, val
  \item \texttt{addEdge}: temp1, \_\_temp3, next
\end{itemize}

\subsection{手動で復元したコード}

\subsubsection{\texttt{lst.append(0)}}
\begin{lstlisting}[language=JavaScript]
var temp1 = new Node();
var temp2 = 0;
temp1.val = temp2;
this.next = temp1;
\end{lstlisting}

\subsubsection{\texttt{lst.append(3)}}
\begin{lstlisting}[language=JavaScript]
var temp1 = new Node();
var temp2 = 3;
temp1.val = temp2;
this.next.next = temp1;
\end{lstlisting}

\subsection{一般化と部分関数の推定}

\begin{lstlisting}[language=JavaScript]
var temp1 = new Node();
var temp2 = f();    // f() = 引数 arg を返す関数
temp1.val = temp2;
g().next = temp1;   // g() = 現在のリストの末尾を返す関数
\end{lstlisting}

\subsection{想定の合成された関数}
\begin{lstlisting}
append(arg) {
  var temp1 = new Node();
  function f() {
    return arg;
  }
  var temp2 = f();
  temp1.val = temp2;
  function g() {
    if(!this.next) {
      return this;
    } else {
      this.next.g();    // できないけど
    }
  }
  g().next = temp1;
}
\end{lstlisting}

\subsection{oracleの関数}
\begin{lstlisting}[language=JavaScript]
append(arg) {
  if(!this.next) {
    this.next = new Node();
    this.net.val = arg;
  } else {
    this.next.append(arg);
  }
}
\end{lstlisting}


% \clearpage

%====================================================================================
\section{対象関数: \texttt{prepend}}
現在のノードを先頭とする連結リストの先頭に、新しいノードを追加するメソッド。
引数として渡された値を \texttt{val} フィールドにもつノードを新たに作成し、それをリストの先頭に接続。

\subsection{初期プログラム}

\begin{lstlisting}[language=JavaScript]
class Node {
  prepend(arg) { /* TBD */ }
}
var lst = new Node(); 
lst.val = 2;
lst = lst.prepend(0);
lst = lst.prepend(3);
\end{lstlisting}

\subsection{操作ログ}

\subsubsection{\texttt{lst.prepend(0)} に対する操作}
\begin{itemize}
  \item \texttt{addNode}: \_\_temp1, Node, false, undefined
  \item \texttt{addNode}: \_\_temp2, 0, true, string
  \item \texttt{addEdge}: \_\_temp1, \_\_temp2, val
  \item \texttt{addEdge}: \_\_temp1, main-new1, next
  \item \texttt{addVariable}: \_\_temp1, return
\end{itemize}

\subsubsection{\texttt{lst.prepend(3)} に対する操作}
\begin{itemize}
  \item \texttt{addNode}: \_\_temp3, Node, false, undefined
  \item \texttt{addNode}: \_\_temp4, 3, true, string
  \item \texttt{addEdge}: \_\_temp3, \_\_temp4, val
  \item \texttt{addEdge}: \_\_temp3, \_\_temp1, next
  \item \texttt{addVariable}: \_\_temp3, return
\end{itemize}

\subsection{手動で復元したコード}

\subsubsection{\texttt{lst.prepend(0)}}
\begin{lstlisting}[language=JavaScript]
var temp1 = new Node();
var temp2 = 0;
temp1.val = temp2;
temp1.next = this;
return temp1;
\end{lstlisting}

\subsubsection{\texttt{lst.prepend(3)}}
\begin{lstlisting}[language=JavaScript]
var temp3 = new Node();
var temp4 = 3;
temp3.val = temp4;
temp3.next = this;
return temp3;
\end{lstlisting}

temp1をthisにすることは実行時情報から判断できる？

\subsection{一般化と部分関数の推定}

\begin{lstlisting}[language=JavaScript]
var temp1 = new Node();
var temp2 = f();    // f() = 引数 arg を返す関数
temp1.val = temp2;
temp1.next = this;
return temp1;
\end{lstlisting}

\subsection{想定の合成された関数}
\begin{lstlisting}
prepend(arg) {
  var temp1 = new Node();
  function f() {
    return arg;
  }
  var temp2 = f();
  temp1.val = temp2;
  temp1.next = this;
  return temp1;
}
\end{lstlisting}

\subsection{oracleの関数}
\begin{lstlisting}[language=JavaScript]
prepend(arg) {
  var temp1 = new Node();
  temp1.val = arg;
  temp1.next = this;
  return temp1;
}
\end{lstlisting}


% \clearpage


%====================================================================================
\section{対象関数: \texttt{removeLast}}
現在のノードを先頭とする連結リストの最後尾のノードを消すメソッド。

\subsection{初期プログラム}

\begin{lstlisting}[language=JavaScript]
class Node {
    removeLast() { /* TBD */ }
}
var lst = new Node(); 
lst.val = 2;
lst.next = new Node();
lst.next.val = 0;
lst.next.next = new Node();
lst.next.next.val = 3;
lst.removeLast();
lst.removeLast();
\end{lstlisting}

\subsection{操作ログ}

\subsubsection{\texttt{lst.removeLast()} に対する操作}
\begin{itemize}
  \item \texttt{deleteNode}: main-new3
\end{itemize}

\subsubsection{\texttt{lst.removeLasr()} に対する操作}
\begin{itemize}
  \item \texttt{deleteNode}: main-new2
\end{itemize}

nextとかvalとかも消しても良い  \\
Nodeだけ消せば可視化されない

\subsection{手動で復元したコード}

\subsubsection{\texttt{lst.removeLast()}}
\begin{lstlisting}[language=JavaScript]
this.next.next = undefined
\end{lstlisting}

\subsubsection{\texttt{lst.removeLast()}}
\begin{lstlisting}[language=JavaScript]
this.next = undefined
\end{lstlisting}

next先のNodeが消えた  \\
nullではなく、undefined  \\
constructorがあるならnullでよい 

\subsection{一般化と部分関数の推定}

\begin{lstlisting}[language=JavaScript]
f().next = undefined;   // f() = 現在のリストの末尾を返す関数
\end{lstlisting}

\subsection{想定の合成された関数}
\begin{lstlisting}
removeLast() {
  function f() {
    if(!this.next.next) {
      return this.next;
    } else {
      this.next.g();    // できないけど
    }
  }
  f().next = undefined;
}
\end{lstlisting}

\subsection{oracleの関数}
\begin{lstlisting}[language=JavaScript]
removeLast() {
  if(this.next.next === undefined) {
    this.next = undefined;
  } else {
    this.next.removeLast();
  }
}
\end{lstlisting}

% \clearpage





%====================================================================================
\section{対象関数: \texttt{removeFirst}}
現在のノードを先頭とする連結リストの先頭のノードを消すメソッド。

\subsection{初期プログラム}

\begin{lstlisting}[language=JavaScript]
class Node {
    removeFirst() { /* TBD */ }
}
var lst = new Node(); 
lst.val = 2;
lst.next = new Node();
lst.next.val = 0;
lst.next.next = new Node();
lst.next.next.val = 3;
lst = lst.removeFirst();
lst = lst.removeFirst();
\end{lstlisting}

\subsection{操作ログ}

\subsubsection{\texttt{lst.removeFirst()} に対する操作}
\begin{itemize}
  \item \texttt{addVariable}: main-new2, return
\end{itemize}

\subsubsection{\texttt{lst.removeFirst()} に対する操作}
\begin{itemize}
  \item \texttt{addVariable}: main-new3, return
\end{itemize}

\subsection{手動で復元したコード}

\subsubsection{\texttt{lst.prepend(0)}}
\begin{lstlisting}[language=JavaScript]
return this.next;
\end{lstlisting}
ここの変換は非自明？


\subsubsection{\texttt{lst.prepend(3)}}
\begin{lstlisting}[language=JavaScript]
return this.next;
\end{lstlisting}

\subsection{一般化と部分関数の推定}

\begin{lstlisting}[language=JavaScript]
return this.next;
\end{lstlisting}

\subsection{想定の合成された関数}
\begin{lstlisting}
removeFirst() {
  return this.next;
}
\end{lstlisting}

\subsection{oracleの関数}
\begin{lstlisting}[language=JavaScript]
removeFirst() {
  return this.next;
}
\end{lstlisting}

% \clearpage





%====================================================================================



%====================================================================================
\section{対象関数: \texttt{removeVal(arg)}}
連結リスト内のvalフィールドがargの値を持つノードを削除するメソッド。

\subsection{初期プログラム}

\begin{lstlisting}[language=JavaScript]
class Node {
  removeVal(arg) { /* TBD */ }
}
var lst = new Node(); 
lst.val = 2;
lst.next = new Node();
lst.next.val = 0;
lst.next.next = new Node(); 
lst.next.next.val = 3;
lst.next.next.next = new Node(); 
lst.next.next.next.val = 1;
lst = lst.removeVal(0);
lst = lst.removeVal(1);
\end{lstlisting}

\subsection{操作ログ}

\subsubsection{\texttt{lst.removeVal(0)} に対する操作}
\begin{itemize}
  \item \texttt{editEdgeReference}: main-new1, main-new2, main-new3, next
  \item \texttt{addVariable}: main-new1, return
\end{itemize}

\subsubsection{\texttt{lst.removeVal(1)} に対する操作}
\begin{itemize}
  \item \texttt{deleteEdge}: main-new3, main-new4, next
  \item \texttt{addVariable}: main-new1, return
\end{itemize}

\subsection{手動で復元したコード}

\subsubsection{\texttt{lst.removeVal(0)}}
\begin{lstlisting}[language=JavaScript]
this.next = this.next.next;
return this;
\end{lstlisting}

\subsubsection{\texttt{lst.removeVal(1)}}
\begin{lstlisting}[language=JavaScript]
this.next.next = null;
return this;
\end{lstlisting}

\subsection{一般化と部分関数の推定}

\begin{lstlisting}[language=JavaScript]
f().next = g()
// f は val フィールドに arg をもつノードの一個前
// g は val フィールドに arg をもつノードの一個後
\end{lstlisting}

\subsection{想定の合成された関数}
\begin{lstlisting}
removeVal(arg) {
  f().next = g();
  return this;
  f() {
    if(this.val === arg) {
      
    }
  }
  g() {
    if(this.val === arg) {
    
    }
  }
}
\end{lstlisting}

\subsection{oracleの関数}
\begin{lstlisting}[language=JavaScript]
removeVal(arg) {
  if(this.val === arg) {
    return this.next;
  }
  if(this.next !== null) {
    this.next = this.next.removeval(arg);
  }
  return this;
}
\end{lstlisting}


% \clearpage




%====================================================================================

%====================================================================================

%====================================================================================

%====================================================================================

%====================================================================================

%====================================================================================

%====================================================================================

%====================================================================================

%====================================================================================

%====================================================================================

%====================================================================================

%====================================================================================

%====================================================================================


%====================================================================================





\section{対象関数: \texttt{removeAt(i)}}
連結リストのi番目のノードを削除するメソッド。

\subsection{初期プログラム}

\begin{lstlisting}[language=JavaScript]
class Node {
  removeAt(i) { /* TBD */ }
}
var lst = new Node(); 
lst.val = 2;
lst.next = new Node();
lst.next.val = 0;
lst.next.next = new Node(); 
lst.next.next.val = 3;
lst.next.next.next = new Node(); 
lst.next.next.next.val = 1;
lst.removeAt(2);
lst.removeAt(1);
\end{lstlisting}

\subsection{操作ログ}

\subsubsection{\texttt{lst.removeAt(2)} に対する操作}
\begin{itemize}
  \item \texttt{editEdgeReference}: main-new2, main-new3, main-new4, next
  \item \texttt{addVariable}: main-new1, return
\end{itemize}

\subsubsection{\texttt{lst.removeAt(1)} に対する操作}
\begin{itemize}
  \item \texttt{editEdgeReference}: main-new1, main-new2, main-new4, next
  \item \texttt{addVariable}: main-new1, return
\end{itemize}

\subsection{手動で復元したコード}

\subsubsection{\texttt{lst.append(0)}}
\begin{lstlisting}[language=JavaScript]
this.next.next = this.next.next.next;
return this;
\end{lstlisting}

\subsubsection{\texttt{lst.append(3)}}
\begin{lstlisting}[language=JavaScript]
this.next = this.next.next;
return this;
\end{lstlisting}

\subsection{一般化と部分関数の推定}

\begin{lstlisting}[language=JavaScript]
f().next = g(); 
\\ f() = i-1 番目のノードを返す関数
\\ g() = i+1 番目のノードを返す関数
return this;
\end{lstlisting}

\subsection{想定の合成された関数}
\begin{lstlisting}
removeAt(i) {
  f().next = g();
  return this;
  \\ f() = i-1 番目のノードを返す関数
  \\ g() = i+1 番目のノードを返す関数
}
\end{lstlisting}

\subsection{oracleの関数}
\begin{lstlisting}[language=JavaScript]
removeAt(i) {
  if(i=== 0) {
    return this.next;
  }
  this.next = this.next.remove(i - 1);
  return this;
}
\end{lstlisting}


% \clearpage






\section{対象関数: \texttt{insertAfter(i, arg)}}
連結リストのi番目の後ろに、valフィールドにargを持つ新しいノードを追加するメソッド。


\subsection{初期プログラム}

\begin{lstlisting}[language=JavaScript]
class Node {
  insertAfter(i, arg) { /* TBD */ }
}
var lst = new Node(); 
lst.val = 2;
lst.next = new Node();
lst.next.val = 0;
lst.insertAfter(0,3);
lst.insertAfter(1,4);
\end{lstlisting}

\subsection{操作ログ}

\subsubsection{\texttt{lst.insertAfter(0,3)} に対する操作}
\begin{itemize}
  \item \texttt{addNode}: \_\_temp1, Node, false, undefined
  \item \texttt{addNode}: \_\_temp2, 3, true, string
  \item \texttt{addEdge}: \_\_temp1, \_\_temp2, val
  \item \texttt{addEdge}: \_\_temp1, main-new2, next
  \item \texttt{editEdgeReference}: main-new1, main-new2, \_\_temp1, next
\end{itemize}

\subsubsection{\texttt{lst.insertAftere(1,4)} に対する操作}
\begin{itemize}
  \item \texttt{addNode}: \_\_temp3, Node, false, undefined
  \item \texttt{addNode}: \_\_temp4, 4, true, string
  \item \texttt{addEdge}: \_\_temp3, \_\_temp4, val
  \item \texttt{addEdge}: \_\_temp3, main-new2, next
  \item \texttt{editEdgeReference}: \_\_temp1, main-new2, \_\_temp3, next
\end{itemize}

\subsection{手動で復元したコード}

\subsubsection{\texttt{lst.insertAfter(0,3)}}
\begin{lstlisting}[language=JavaScript]
var temp1 = new Node();
var temp2 = 3;
temp1.val = temp2;
temp1.next = this.next;
this.next = temp1;
\end{lstlisting}

\subsubsection{\texttt{lst.insetAfter(1,4)}}
\begin{lstlisting}[language=JavaScript]
var temp1 = new Node();
var temp2 = 4;
temp1.val = temp2;
temp1.next = this.next.next;
this.next.next = temp1;
\end{lstlisting}

\subsection{一般化と部分関数の推定}

\begin{lstlisting}[language=JavaScript]
var temp1 = new Node();
var temp2 = f();        // f() = 引数 arg を返す関数
temp1.val = temp2;
temp1.next = g()    // g() = i+1 番目のノードを返す関数
h().next = temp1;   // h() = i-1 番目のノードを返す関数
\end{lstlisting}

\subsection{想定の合成された関数}
\begin{lstlisting}
insertAfter(i,arg) {
  var temp1 = new Node();
  var temp2 = arg;
  temp1.val = temp2;
  temp1.next = find(i+1);
  find(i-1).next = temp1;
}
\end{lstlisting}

\subsection{oracleの関数}
\begin{lstlisting}[language=JavaScript]
insertAfter(i, arg) {
  function find(node, index) {
    if (node === null) {
      return null;
    }
    if (index === 0) {
      return node;
    }
    return find(node.next, index - 1);
  }
  const targetNode = find(this, i);
  if (targetNode !== null) {
    const newNode = new Node();
    newNode.val = arg;
    newNode.next = targetNode.next;
    targetNode.next = newNode;
  }
}
\end{lstlisting}

% \clearpage




\section{対象関数: \texttt{concat(lst)}}
現在のノードを先頭とする連結リストの末尾に、別の連結リストを追加するメソッド。

\subsection{初期プログラム}

\begin{lstlisting}[language=JavaScript]
class Node {
  concat(otherList) { /* TBD */ }
}
var lst = new Node(); 
lst.val = 2;
lst.next = new Node();
lst.next.val = 0;
lst.next.next = new Node(); 
lst.next.next.val = 3;
var list = new Node(); 
list.val = 1;
list.next = new Node();
list.next.val = 4;
var l = new Node();
l.val = 5;
lst.concat(list);
lst.concat(l);
\end{lstlisting}

\subsection{操作ログ}

\subsubsection{\texttt{lst.concat(list)} に対する操作}
\begin{itemize}
  \item \texttt{addEdge}: main-new3, main-new4, next
\end{itemize}

\subsubsection{\texttt{lst.concat(l)} に対する操作}
\begin{itemize}
  \item \texttt{addNode}: main-new5, main-new6, next
\end{itemize}

\subsection{手動で復元したコード}

\subsubsection{\texttt{lst.concat(list)}}
\begin{lstlisting}[language=JavaScript]
this.next.next = list;
\end{lstlisting}

\subsubsection{\texttt{lst.concat(l)}}
\begin{lstlisting}[language=JavaScript]
this.next.next.next.next.next = l;
\end{lstlisting}

\subsection{一般化と部分関数の推定}

\begin{lstlisting}[language=JavaScript]
f().next = g();
\\ f() = 最後のノードを返す関数
\\ g() = 引数を返す関数
\end{lstlisting}

\subsection{想定の合成された関数}
\begin{lstlisting}
concat(otherList) {
  f().next = otherList;
  \\ f() = thisの最後のノードを返す関数
}
\end{lstlisting}

\subsection{oracleの関数}
\begin{lstlisting}[language=JavaScript]
concat(otherList) {
  if(!this.next) {
    this.next = otherList;
  } else {
    this.next.concat(otherList);
  }
}
\end{lstlisting}


% \clearpage




\section{対象関数: \texttt{set(i,arg)}}
連結リストのi番目の値をargに変更するメソッド。

\subsection{初期プログラム}

\begin{lstlisting}[language=JavaScript]
class Node {
  set(i,arg) { /* TBD */ }
}
var lst = new Node(); 
lst.val = 2;
lst.next = new Node();
lst.next.val = 0;
lst.next.next = new Node(); 
lst.next.next.val = 3;
lst.next.next.next = new Node(); 
lst.next.next.next.val = 1;
lst.set(2,5);
lst.set(1,4);
\end{lstlisting}

\subsection{操作ログ}

\subsubsection{\texttt{lst.set(2,5)} に対する操作}
\begin{itemize}
  \item \texttt{addNode}: \_\_temp1, 5, true, string
  \item \texttt{editEdgeReference}: main-new3, main-new3-val, \_\_temp1, val
\end{itemize}

\subsubsection{\texttt{lst.set(1,4)} に対する操作}
\begin{itemize}
  \item \texttt{addNode}: \_\_temp2, 4, true, string
  \item \texttt{addEdge}: main-new2, main-new2-val, \_\_temp2, val
\end{itemize}

\subsection{手動で復元したコード}

\subsubsection{\texttt{lst.set(2,5)}}
\begin{lstlisting}[language=JavaScript]
var temp1 = 5;
this.next.next.val = temp1;
\end{lstlisting}

\subsubsection{\texttt{lst.set(1,4)}}
\begin{lstlisting}[language=JavaScript]
var temp1 = 4;
this.next.val = temp2;
\end{lstlisting}

\subsection{一般化と部分関数の推定}

\begin{lstlisting}[language=JavaScript]
var temp1 = f();    \\ f() = 引数 arg を返す関数
g().val = temp1;    \\ g() = i番目のノードを返す関数
\end{lstlisting}

\subsection{想定の合成された関数}
\begin{lstlisting}
set(i,arg) {
  var temp1 = arg;
  find(i).val = temp1;
}
\end{lstlisting}

\subsection{oracleの関数}
\begin{lstlisting}[language=JavaScript]
set(i,arg) {
  if(i === 0) {
    this.val = arg;
  }
  if(i > 0) {
    this.next.set(i-1,arg);
  }
}
\end{lstlisting}


% \clearpage


\section{対象関数: \texttt{swap(i, j)}}
現在のノードを先頭とする連結リストのi番目とj番目の要素を交換するメソッド。

\subsection{初期プログラム}

\begin{lstlisting}[language=JavaScript]
class Node {
  swap(i,j) { /* TBD */ }
}
var lst = new Node(); 
lst.val = 2;
lst.next = new Node();
lst.next.val = 0;
lst.next.next = new Node(); 
lst.next.next.val = 3;
lst.next.next.next = new Node(); 
lst.next.next.next.val = 1;
lst.swap(0, 3);
lst.swap(1,2);
\end{lstlisting}

\subsection{操作ログ}

\subsubsection{\texttt{lst.swap(0,3)} に対する操作}
\begin{itemize}
  \item \texttt{addNode}: \_\_temp1, 2, true, string
  \item \texttt{addNode}: \_\_temp2, 1, true, string
  \item \texttt{editEdgeReference}: main-new1, main-new1-val, \_\_temp1, val
  \item \texttt{editEdgeReference}: main-new4, main-new4-val, \_\_temp2, val
\end{itemize}

\subsubsection{\texttt{lst.swap(1,2)} に対する操作}
\begin{itemize}
  \item \texttt{addNode}: \_\_temp3, 3, true, string
  \item \texttt{addNode}: \_\_temp4, 0, true, string
  \item \texttt{editEdgeReference}: main-new2, main-new2-val, \_\_temp3, val
  \item \texttt{editEdgeReference}: main-new3, main-new3-val, \_\_temp4, val
\end{itemize}

\subsection{手動で復元したコード}
操作列だけみると実際には違う挙動をするプログラムが出てくる   \\
参照元がなくなったが使うものは先にtempとかで変数宣言しておく？
\subsubsection{\texttt{lst.swap(0,3)}}
\begin{lstlisting}[language=JavaScript]
var temp1 = 2;
var temp2 = 1;
this.val = temp1;
this.next.next.next.val = temp2;
\end{lstlisting}

\subsubsection{\texttt{lst.swap(1,2)}}
\begin{lstlisting}[language=JavaScript]
var temp1 = 3;
var temp2 = 0;
this.next.val = temp1;
this.next.next.val = temp2;
\end{lstlisting}

\subsection{一般化と部分関数の推定}
\begin{lstlisting}[language=JavaScript]
var temp1 = f();    \\ f() = i 番目のノードの val フィールドの値を返す関数
var temp2 = g();    \\ g() = j 番目のノードの val フィールドの値を返す関数
h().val = temp1;    \\ h() = j 番目のノードを返す関数
i().val = temp2;    \\ i() = i 番目のノードを返す関数
\end{lstlisting}

\subsection{想定の合成された関数}
\begin{lstlisting}
swap(i,j) {
  var temp1 = getVal(i);
  var temp2 = getVal(j);
  getNode(j).val = temp1;
  getNode(i).val = temp2;
}
\end{lstlisting}

\subsection{oracleの関数}
\begin{lstlisting}[language=JavaScript]
swap(i,j) {
  function find(node, index) {
    if(index === 0) {
      return node;
    }
    return find(node.next, index - 1);
  }
  var nodeI = find(this, i);
  var nodeJ = find(this, j);
  var tempVal = nodeI.val;
  nodeI.val = nodeJ.val;
  nodeJ.val = tempVal;
}
\end{lstlisting}


% \clearpage

\section{対象関数: \texttt{indexOf(arg)}}
連結リスト内で指定された値を持つノードのインデックスを返すメソッド。

\subsection{初期プログラム}

\begin{lstlisting}[language=JavaScript]
class Node {
 indexOf(arg) { /* TBD */ }
}
var lst = new Node(); 
lst.val = 2;
lst.next = new Node();
lst.next.val = 0;
lst.next.next = new Node();
lst.next.next.val = 3;
var index0 = lst.indexOf(0);
var index3 = lst.indexOf(3);
var indexNotFound = lst.indexOf(99);
\end{lstlisting}

\subsection{操作ログ}

\subsubsection{\texttt{lst.indexOf(0)} に対する操作}
\begin{itemize}
  \item \texttt{addVariable}: main-new2, return
\end{itemize}

\subsubsection{\texttt{lst.indexOf(3)} に対する操作}
\begin{itemize}
  \item \texttt{addVariable}: main-new3, return
\end{itemize}

\subsubsection{\texttt{lst.indexOf(99)} に対する操作}
\begin{itemize}
  \item \texttt{addVariable}: null, return
\end{itemize}

\subsection{手動で復元したコード}

\subsubsection{\texttt{lst.indexOf(0)}}
\begin{lstlisting}[language=JavaScript]
if (this.val === 0) return 0;
if (this.next.val === 0) return 1;
// ...
\end{lstlisting}

\subsubsection{\texttt{lst.indexOf(3)}}
\begin{lstlisting}[language=JavaScript]
if (this.val === 3) return 0;
if (this.next.val === 3) return 1;
if (this.next.next.val === 3) return 2;
// ...
\end{lstlisting}

\subsection{一般化と部分関数の推定}

\begin{lstlisting}[language=JavaScript]
// f() = 現在のノードからargを探し、インデックスを返す関数
return f();
\end{lstlisting}

\subsection{想定の合成された関数}
\begin{lstlisting}
indexOf(arg) {
  let currentIndex = 0;
  let currentNode = this;
  while (currentNode !== null) {
    if (currentNode.val === arg) {
      return currentIndex;
    }
    currentNode = currentNode.next;
    currentIndex++;
  }
  return -1; // Not found
}
\end{lstlisting}

\subsection{oracleの関数}
\begin{lstlisting}[language=JavaScript]
indexOf(arg, current=this, index=0) {
  if (current === null) {
    return -1;
  }
  if (current.val === arg) {
    return index;
  }
  return this.indexOf(arg, current.next, index + 1);
}
\end{lstlisting}

% \clearpage

\section{対象関数: \texttt{size()}}
連結リストのノード数を返すメソッド。

\subsection{初期プログラム}

\begin{lstlisting}[language=JavaScript]
class Node {
  size() { /* TBD */ }
}
var lst = new Node(); 
lst.val = 2;
lst.next = new Node();
lst.next.val = 0;
lst.next.next = new Node();
lst.next.next.val = 3;
var s1 = lst.size();
var emptyLst = new Node();
var s2 = emptyLst.size();
\end{lstlisting}

\subsection{操作ログ}

\subsubsection{\texttt{lst.size()} に対する操作}
\begin{itemize}
  \item \texttt{addVariable}: \_\_temp1, 3, true, number
  \item \texttt{addVariable}: \_\_temp1, return
\end{itemize}

\subsubsection{\texttt{emptyLst.size()} に対する操作}
\begin{itemize}
  \item \texttt{addVariable}: \_\_temp2, 1, true, number
  \item \texttt{addVariable}: \_\_temp2, return
\end{itemize}

\subsection{手動で復元したコード}

\subsubsection{\texttt{lst.size()}}
\begin{lstlisting}[language=JavaScript]
var count = 0;
var current = this;
while (current !== null) {
  count++;
  current = current.next;
}
return count;
\end{lstlisting}

\subsubsection{\texttt{emptyLst.size()}}
\begin{lstlisting}[language=JavaScript]
  return 1; // Assuming initial node counts as 1 even if val is undefined
\end{lstlisting}

\subsection{一般化と部分関数の推定}

\begin{lstlisting}[language=JavaScript]
// f() = リストを走査してノード数を数える関数
return f();
\end{lstlisting}

\subsection{想定の合成された関数}
\begin{lstlisting}
size() {
  let count = 0;
  let current = this;
  while (current !== null) {
    count++;
    current = current.next;
  }
  return count;
}
\end{lstlisting}

\subsection{oracleの関数}
\begin{lstlisting}[language=JavaScript]
size(node = this) {
  if (node === null) {
    return 0;
  }
  return 1 + this.size(node.next);
}
\end{lstlisting}

% \clearpage

\section{対象関数: \texttt{get(i)}}
連結リストのi番目のノードの\texttt{val}フィールドを返すメソッド。

\subsection{初期プログラム}

\begin{lstlisting}[language=JavaScript]
class Node {
  get(i) { /* TBD */ }
}
var lst = new Node(); 
lst.val = 2;
lst.next = new Node();
lst.next.val = 0;
lst.next.next = new Node();
lst.next.next.val = 3;
var val0 = lst.get(0);
var val2 = lst.get(2);
\end{lstlisting}

\subsection{操作ログ}

\subsubsection{\texttt{lst.get(0)} に対する操作}
\begin{itemize}
  \item \texttt{addVariable}: main-new1-val, return
\end{itemize}

\subsubsection{\texttt{lst.get(2)} に対する操作}
\begin{itemize}
  \item \texttt{addVariable}: main-new3-val, return
\end{itemize}

\subsection{手動で復元したコード}

\subsubsection{\texttt{lst.get(0)}}
\begin{lstlisting}[language=JavaScript]
return this.val;
\end{lstlisting}

\subsubsection{\texttt{lst.get(2)}}
\begin{lstlisting}[language=JavaScript]
return this.next.next.val;
\end{lstlisting}

\subsection{一般化と部分関数の推定}

\begin{lstlisting}[language=JavaScript]
// f() = i番目のノードのvalフィールドを返す関数
return f();
\end{lstlisting}

\subsection{想定の合成された関数}
\begin{lstlisting}
get(i) {
  let currentIndex = 0;
  let currentNode = this;
  while (currentNode !== null && currentIndex < i) {
    currentNode = currentNode.next;
    currentIndex++;
  }
  if (currentNode !== null) {
    return currentNode.val;
  }
  return undefined; // Or throw error
}
\end{lstlisting}

\subsection{oracleの関数}
\begin{lstlisting}[language=JavaScript]
get(i, node = this) {
  if (node === null) {
    return undefined; // Or throw error
  }
  if (i === 0) {
    return node.val;
  }
  return this.get(i - 1, node.next);
}
\end{lstlisting}

% \clearpage

\section{対象関数: \texttt{clear()}}
連結リストのすべてのノードを削除するメソッド。

\subsection{初期プログラム}

\begin{lstlisting}[language=JavaScript]
class Node {
  clear() { /* TBD */ }
}
var lst = new Node(); 
lst.val = 2;
lst.next = new Node();
lst.next.val = 0;
lst.next.next = new Node();
lst.next.next.val = 3;
lst.clear();
var emptyLst = new Node();
emptyLst.clear();
\end{lstlisting}

\subsection{操作ログ}

\subsubsection{\texttt{lst.clear()} に対する操作}
\begin{itemize}
  \item \texttt{deleteNode}: main-new3
  \item \texttt{deleteNode}: main-new2
  \item \texttt{deleteNode}: main-new1
\end{itemize}

\subsubsection{\texttt{emptyLst.clear()} に対する操作}
\begin{itemize}
  \item (No explicit deleteNode operations if the initial node is considered the list itself and its next is null)
\end{itemize}

\subsection{手動で復元したコード}

\subsubsection{\texttt{lst.clear()}}
\begin{lstlisting}[language=JavaScript]
this.val = undefined; // Or some default value
this.next = undefined;
\end{lstlisting}

\subsubsection{\texttt{emptyLst.clear()}}
\begin{lstlisting}[language=JavaScript]
// No change if already effectively empty
\end{lstlisting}

\subsection{一般化と部分関数の推定}

\begin{lstlisting}[language=JavaScript]
// f() = リストを走査してノードを削除する関数
f();
\end{lstlisting}

\subsection{想定の合成された関数}
\begin{lstlisting}
clear() {
  // In JavaScript, setting next to null effectively "clears" the list
  // if you assume the current node remains the head.
  // If the head itself should be "emptied", then its val should be reset.
  if (this.next) {
    // Recursively clear the rest of the list
    this.next.clear();
    // Then remove the reference to the cleared tail
    this.next = undefined;
  }
  // Reset the current node's value
  this.val = undefined;
}
\end{lstlisting}

\subsection{oracleの関数}
\begin{lstlisting}[language=JavaScript]
clear() {
  if (this.next !== null) {
    this.next.clear();
    this.next = null;
  }
  this.val = undefined; // Or some default value
}
\end{lstlisting}

% \clearpage

\section{対象関数: \texttt{contains(arg)}}
連結リスト内に指定された値を持つノードが存在するかどうかを返すメソッド。

\subsection{初期プログラム}

\begin{lstlisting}[language=JavaScript]
class Node {
  contains(arg) { /* TBD */ }
}
var lst = new Node(); 
lst.val = 2;
lst.next = new Node();
lst.next.val = 0;
lst.next.next = new Node();
lst.next.next.val = 3;
var c1 = lst.contains(0);
var c2 = lst.contains(99);
\end{lstlisting}

\subsection{操作ログ}

\subsubsection{\texttt{lst.contains(0)} に対する操作}
\begin{itemize}
  \item \texttt{addVariable}: \_\_temp1, true, true, boolean
  \item \texttt{addVariable}: \_\_temp1, return
\end{itemize}

\subsubsection{\texttt{lst.contains(99)} に対する操作}
\begin{itemize}
  \item \texttt{addVariable}: \_\_temp2, false, true, boolean
  \item \texttt{addVariable}: \_\_temp2, return
\end{itemize}

\subsection{手動で復元したコード}

\subsubsection{\texttt{lst.contains(0)}}
\begin{lstlisting}[language=JavaScript]
if (this.val === 0) return true;
if (this.next && this.next.val === 0) return true;
return false;
\end{lstlisting}

\subsubsection{\texttt{lst.contains(99)}}
\begin{lstlisting}[language=JavaScript]
var current = this;
while (current !== null) {
  if (current.val === 99) return true;
  current = current.next;
}
return false;
\end{lstlisting}

\subsection{一般化と部分関数の推定}

\begin{lstlisting}[language=JavaScript]
// f() = リストを走査してargを探し、見つかればtrue、そうでなければfalseを返す関数
return f();
\end{lstlisting}

\subsection{想定の合成された関数}
\begin{lstlisting}
contains(arg) {
  let current = this;
  while (current !== null) {
    if (current.val === arg) {
      return true;
    }
    current = current.next;
  }
  return false;
}
\end{lstlisting}

\subsection{oracleの関数}
\begin{lstlisting}[language=JavaScript]
contains(arg) {
  if (this === null) {
    return false;
  }
  if (this.val === arg) {
    return true;
  }
  if (this.next !== null) {
    return this.next.contains(arg);
  }
  return false;
}
\end{lstlisting}

% \clearpage

\section{対象関数: \texttt{isEmpty()}}
連結リストが空かどうかを返すメソッド。

\subsection{初期プログラム}

\begin{lstlisting}[language=JavaScript]
class Node {
  isEmpty() { /* TBD */ }
}
var lst = new Node(); 
lst.val = 2;
lst.next = new Node();
lst.next.val = 0;
var e1 = lst.isEmpty();
var emptyLst = new Node();
var e2 = emptyLst.isEmpty();
\end{lstlisting}

\subsection{操作ログ}

\subsubsection{\texttt{lst.isEmpty()} に対する操作}
\begin{itemize}
  \item \texttt{addVariable}: \_\_temp1, false, true, boolean
  \item \texttt{addVariable}: \_\_temp1, return
\end{itemize}

\subsubsection{\texttt{emptyLst.isEmpty()} に対する操作}
\begin{itemize}
  \item \texttt{addVariable}: \_\_temp2, true, true, boolean
  \item \texttt{addVariable}: \_\_temp2, return
\end{itemize}

\subsection{手動で復元したコード}

\subsubsection{\texttt{lst.isEmpty()}}
\begin{lstlisting}[language=JavaScript]
return this.val === undefined && this.next === undefined; // Or check for specific initial state
\end{lstlisting}

\subsubsection{\texttt{emptyLst.isEmpty()}}
\begin{lstlisting}[language=JavaScript]
return true;
\end{lstlisting}

\subsection{一般化と部分関数の推定}

\begin{lstlisting}[language=JavaScript]
// f() = リストが空かどうかを判定する関数
return f();
\end{lstlisting}

\subsection{想定の合成された関数}
\begin{lstlisting}
isEmpty() {
  // Assuming an "empty" list has a null 'next' and potentially an undefined 'val'
  return this.next === null && this.val === undefined;
}
\end{lstlisting}

\subsection{oracleの関数}
\begin{lstlisting}[language=JavaScript]
isEmpty() {
  return this.next === null && this.val === undefined;
}
\end{lstlisting}

% \clearpage

\section{対象関数: \texttt{update(oldVal, newVal)}}
連結リスト内で\texttt{oldVal}を持つすべてのノードの値を\texttt{newVal}に変更するメソッド。

\subsection{初期プログラム}

\begin{lstlisting}[language=JavaScript]
class Node {
  update(oldVal, newVal) { /* TBD */ }
}
var lst = new Node(); 
lst.val = 2;
lst.next = new Node();
lst.next.val = 0;
lst.next.next = new Node();
lst.next.next.val = 3;
lst.next.next.next = new Node();
lst.next.next.next.val = 0; // Another 0
lst.update(0, 99);
lst.update(5, 100); // Not found
\end{lstlisting}

\subsection{操作ログ}

\subsubsection{\texttt{lst.update(0, 99)} に対する操作}
\begin{itemize}
  \item \texttt{addNode}: \_\_temp1, 99, true, string
  \item \texttt{editEdgeReference}: main-new2, main-new2-val, \_\_temp1, val
  \item \texttt{addNode}: \_\_temp2, 99, true, string
  \item \texttt{editEdgeReference}: main-new4, main-new4-val, \_\_temp2, val
\end{itemize}

\subsubsection{\texttt{lst.update(5, 100)} に対する操作}
\begin{itemize}
  \item (No operations if value not found)
\end{itemize}

\subsection{手動で復元したコード}

\subsubsection{\texttt{lst.update(0, 99)}}
\begin{lstlisting}[language=JavaScript]
var current = this;
while (current !== null) {
  if (current.val === 0) {
    current.val = 99;
  }
  current = current.next;
}
\end{lstlisting}

\subsection{一般化と部分関数の推定}

\begin{lstlisting}[language=JavaScript]
// f() = リストを走査してoldValを探し、newValに変更する関数
f();
\end{lstlisting}

\subsection{想定の合成された関数}
\begin{lstlisting}
update(oldVal, newVal) {
  let current = this;
  while (current !== null) {
    if (current.val === oldVal) {
      current.val = newVal; // This implies addNode and editEdgeReference for val
    }
    current = current.next;
  }
}
\end{lstlisting}

\subsection{oracleの関数}
\begin{lstlisting}[language=JavaScript]
update(oldVal, newVal) {
  if (this === null) {
    return;
  }
  if (this.val === oldVal) {
    this.val = newVal;
  }
  if (this.next !== null) {
    this.next.update(oldVal, newVal);
  }
}
\end{lstlisting}

% \clearpage

\section{対象関数: \texttt{filter(predicate)}}
指定された条件に一致するノードのみを含む新しい連結リストを返すメソッド。

\subsection{初期プログラム}

\begin{lstlisting}[language=JavaScript]
class Node {
  filter(predicate) { /* TBD */ }
}
var lst = new Node(); 
lst.val = 2;
lst.next = new Node();
lst.next.val = 0;
lst.next.next = new Node();
lst.next.next.val = 3;
lst.next.next.next = new Node();
lst.next.next.next.val = 1;
var evenList = lst.filter(val => val % 2 === 0);
var oddList = lst.filter(val => val % 2 !== 0);
\end{lstlisting}

\subsection{操作ログ}

\subsubsection{\texttt{lst.filter(val => val \% 2 === 0)} に対する操作}
\begin{itemize}
  \item \texttt{addNode}: filteredListHead, Node, false, undefined
  \item \texttt{addNode}: filteredListHeadVal, 2, true, string
  \item \texttt{addEdge}: filteredListHead, filteredListHeadVal, val
  \item \texttt{addNode}: filteredListSecond, Node, false, undefined
  \item \texttt{addNode}: filteredListSecondVal, 0, true, string
  \item \texttt{addEdge}: filteredListSecond, filteredListSecondVal, val
  \item \texttt{addEdge}: filteredListHead, filteredListSecond, next
  \item \texttt{addVariable}: filteredListHead, return
\end{itemize}

\subsubsection{\texttt{lst.filter(val => val \% 2 !== 0)} に対する操作}
\begin{itemize}
  \item \texttt{addNode}: oddListHead, Node, false, undefined
  \item \texttt{addNode}: oddListHeadVal, 3, true, string
  \item \texttt{addEdge}: oddListHead, oddListHeadVal, val
  \item \texttt{addNode}: oddListSecond, Node, false, undefined
  \item \texttt{addNode}: oddListSecondVal, 1, true, string
  \item \texttt{addEdge}: oddListSecond, oddListSecondVal, val
  \item \texttt{addEdge}: oddListHead, oddListSecond, next
  \item \texttt{addVariable}: oddListHead, return
\end{itemize}

\subsection{手動で復元したコード}

\subsubsection{\texttt{lst.filter(val => val \% 2 === 0)}}
\begin{lstlisting}[language=JavaScript]
var newList = null;
var tail = null;
var current = this;
while (current !== null) {
  if (current.val % 2 === 0) {
    var newNode = new Node();
    newNode.val = current.val;
    if (newList === null) {
      newList = newNode;
      tail = newNode;
    } else {
      tail.next = newNode;
      tail = newNode;
    }
  }
  current = current.next;
}
return newList;
\end{lstlisting}

\subsection{一般化と部分関数の推定}

\begin{lstlisting}[language=JavaScript]
// f() = predicateに基づいて新しいリストを構築する関数
return f();
\end{lstlisting}

\subsection{想定の合成された関数}
\begin{lstlisting}
filter(predicate) {
  let newList = null;
  let currentNewNode = null;
  let currentOriginalNode = this;

while (currentOriginalNode !== null) {
    if (predicate(currentOriginalNode.val)) {
      const newNode = new Node();
      newNode.val = currentOriginalNode.val;
      if (newList === null) {
        newList = newNode;
        currentNewNode = newNode;
      } else {
        currentNewNode.next = newNode;
        currentNewNode = newNode;
      }
    }
    currentOriginalNode = currentOriginalNode.next;
  }
  return newList;
}
\end{lstlisting}

\subsection{oracleの関数}
\begin{lstlisting}[language=JavaScript]
filter(predicate) {
  if (this === null) {
    return null;
  }
  if (predicate(this.val)) {
    const newNode = new Node();
    newNode.val = this.val;
    newNode.next = this.next ? this.next.filter(predicate) : null;
    return newNode;
  } else {
    return this.next ? this.next.filter(predicate) : null;
  }
}
\end{lstlisting}

% \clearpage

\section{対象関数: \texttt{deduplicate()}}
連結リストから重複する値を削除し、一意の値を保持するメソッド。

\subsection{初期プログラム}

\begin{lstlisting}[language=JavaScript]
class Node {
  deduplicate() { /* TBD */ }
}
var lst1 = new Node(); 
lst1.val = 1;
lst1.next = new Node();
lst1.next.val = 2;
lst1.next.next = new Node();
lst1.next.next.val = 1; // Duplicate
lst1.next.next.next = new Node();
lst1.next.next.next.val = 3;
lst1.deduplicate();

var lst2 = new Node();
lst2.val = 1;
lst2.next = new Node();
lst2.next.val = 1;
lst2.next.next = new Node();
lst2.next.next.val = 1; // Multiple duplicates
lst2.deduplicate();
\end{lstlisting}

\subsection{操作ログ}

\subsubsection{\texttt{lst1.deduplicate()} に対する操作}
\begin{itemize}
  \item \texttt{deleteNode}: main-new3
  \item \texttt{editEdgeReference}: main-new2, main-new3, main-new4, next
\end{itemize}

\subsubsection{\texttt{lst2.deduplicate()} に対する操作}
\begin{itemize}
  \item \texttt{deleteNode}: main-new6
  \item \texttt{editEdgeReference}: main-new5, main-new6, main-new7, next
  \item \texttt{deleteNode}: main-new7
  \item \texttt{editEdgeReference}: main-new5, main-new7, main-new8, next
\end{itemize}

\subsection{手動で復元したコード}

\subsubsection{\texttt{lst1.deduplicate()}}
\begin{lstlisting}[language=JavaScript]
var current = this;
var seen = new Set();
if (current) {
  seen.add(current.val);
}
while (current && current.next) {
  if (seen.has(current.next.val)) {
    current.next = current.next.next;
  } else {
    seen.add(current.next.val);
    current = current.next;
  }
}
\end{lstlisting}

\subsection{一般化と部分関数の推定}

\begin{lstlisting}[language=JavaScript]
// f() = リストを走査し、重複を検出し、ノードを削除する関数
f();
\end{lstlisting}

\subsection{想定の合成された関数}
\begin{lstlisting}
deduplicate() {
  let seen = new Set();
  let current = this;
  let prev = null;

while (current !== null) {
    if (seen.has(current.val)) {
      // Delete the current node
      if (prev !== null) {
        prev.next = current.next; // editEdgeReference
      } else {
        // If current is head and is a duplicate, the head shifts
        // This case is complex as 'this' itself would need to change.
        // For simplicity, assuming 'this' is not removed directly,
        // or handled by a return value.
      }
      // Implicit deleteNode for current
    } else {
      seen.add(current.val);
      prev = current;
    }
    current = current.next;
  }
}
\end{lstlisting}

\subsection{oracleの関数}
\begin{lstlisting}[language=JavaScript]
deduplicate() {
  let seen = new Set();
  let current = this;
  let prev = null;

while (current !== null) {
    if (seen.has(current.val)) {
      // Skip current node (effectively deleting it)
      if (prev !== null) {
        prev.next = current.next;
      } else {
        // If head is a duplicate, advance the head (caller responsibility)
        // This requires 'return this.next' logic
      }
    } else {
      seen.add(current.val);
      prev = current;
    }
    current = current.next;
  }
}
\end{lstlisting>

% \clearpage

\section{対象関数: \texttt{reverse()}}
連結リストのノードの順序を反転させるメソッド。

\subsection{初期プログラム}

\begin{lstlisting}[language=JavaScript]
class Node {
  reverse() { /* TBD */ }
}
var lst = new Node(); 
lst.val = 1;
lst.next = new Node();
lst.next.val = 2;
lst.next.next = new Node();
lst.next.next.val = 3;
lst.reverse();
\end{lstlisting}

\subsection{操作ログ}

\subsubsection{\texttt{lst.reverse()} に対する操作}
\begin{itemize}
  \item \texttt{editEdgeReference}: main-new3, null, main-new2, next
  \item \texttt{editEdgeReference}: main-new2, main-new3, main-new1, next
  \item \texttt{editEdgeReference}: main-new1, main-new2, null, next
  \item \texttt{addVariable}: main-new3, return
\end{itemize}

\subsection{手動で復元したコード}

\subsubsection{\texttt{lst.reverse()}}
\begin{lstlisting}[language=JavaScript]
var prev = null;
var current = this;
while (current !== null) {
  var nextTemp = current.next;
  current.next = prev;
  prev = current;
  current = nextTemp;
}
return prev;
\end{lstlisting>

\subsection{一般化と部分関数の推定}

\begin{lstlisting}[language=JavaScript]
// f() = ポインタを操作してリストを反転させる関数
return f();
\end{lstlisting}

\subsection{想定の合成された関数}
\begin{lstlisting}
reverse() {
  let prev = null;
  let current = this;
  let next = null;

while (current !== null) {
    next = current.next; // Store next node
    current.next = prev; // Reverse current node's next pointer (editEdgeReference)
    prev = current; // Move prev to current node
    current = next; // Move current to next node
  }
  // At the end, prev will be the new head
  return prev;
}
\end{lstlisting}

\subsection{oracleの関数}
\begin{lstlisting}[language=JavaScript]
reverse(head = this, prev = null) {
  if (head === null) {
    return prev;
  }
  const nextNode = head.next;
  head.next = prev;
  return this.reverse(nextNode, head);
}
\end{lstlisting}

% \clearpage

\section{対象関数: \texttt{sort()}}
連結リストのノードを値に基づいてソートするメソッド。

\subsection{初期プログラム}

\begin{lstlisting}[language=JavaScript]
class Node {
  sort() { /* TBD */ }
}
var lst = new Node(); 
lst.val = 3;
lst.next = new Node();
lst.next.val = 1;
lst.next.next = new Node();
lst.next.next.val = 2;
lst.sort();
\end{lstlisting}

\subsection{操作ログ}

\subsubsection{\texttt{lst.sort()} に対する操作}
\begin{itemize}
  \item \texttt{addNode}: \_\_temp1, 1, true, string
  \item \texttt{editEdgeReference}: main-new1, main-new1-val, \_\_temp1, val
  \item \texttt{addNode}: \_\_temp2, 2, true, string
  \item \texttt{editEdgeReference}: main-new2, main-new2-val, \_\_temp2, val
  \item \texttt{addNode}: \_\_temp3, 3, true, string
  \item \texttt{editEdgeReference}: main-new3, main-new3-val, \_\_temp3, val
\end{itemize}

\subsection{手動で復元したコード}

\subsubsection{\texttt{lst.sort()}}
\begin{lstlisting}[language=JavaScript]
// Example: Bubble sort on values
var swapped;
do {
  swapped = false;
  var current = this;
  while (current !== null && current.next !== null) {
    if (current.val > current.next.val) {
      var temp = current.val;
      current.val = current.next.val;
      current.next.val = temp;
      swapped = true;
    }
    current = current.next;
  }
} while (swapped);
\end{lstlisting}

\subsection{一般化と部分関数の推定}

\begin{lstlisting}[language=JavaScript]
// f() = ソートアルゴリズムを実装し、ノードの値を交換する関数
f();
\end{lstlisting}

\subsection{想定の合成された関数}
\begin{lstlisting}
sort() {
  // Simple bubble sort (modifying values directly)
  let swapped;
  do {
    swapped = false;
    let current = this;
    while (current !== null && current.next !== null) {
      if (current.val > current.next.val) {
        // Swap values (implies addNode for temp val, editEdgeReference for val)
        let temp = current.val;
        current.val = current.next.val;
        current.next.val = temp;
        swapped = true;
      }
      current = current.next;
    }
  } while (swapped);
}
\end{lstlisting}

\subsection{oracleの関数}
\begin{lstlisting}[language=JavaScript]
sort() {
  // For simplicity, converting to array, sorting, then back to list.
  // A true in-place linked list sort would be more complex and recursive.
  const arr = [];
  let current = this;
  while (current !== null) {
    arr.push(current.val);
    current = current.next;
  }
  arr.sort((a, b) => a - b);

current = this;
  let i = 0;
  while (current !== null) {
    current.val = arr[i++];
    current = current.next;
  }
}
\end{lstlisting}

% \clearpage

\section{対象関数: \texttt{splitAt(i)}}
連結リストをi番目のノードで2つのリストに分割するメソッド。元のリストはi番目のノードまでとなり、i番目以降は新しいリストとなる。

\subsection{初期プログラム}

\begin{lstlisting}[language=JavaScript]
class Node {
  splitAt(i) { /* TBD */ }
}
var lst = new Node(); 
lst.val = 1;
lst.next = new Node();
lst.next.val = 2;
lst.next.next = new Node();
lst.next.next.val = 3;
lst.next.next.next = new Node();
lst.next.next.next.val = 4;
var secondList = lst.splitAt(1);
var thirdList = lst.splitAt(0);
\end{lstlisting}

\subsection{操作ログ}

\subsubsection{\texttt{lst.splitAt(1)} に対する操作}
\begin{itemize}
  \item \texttt{editEdgeReference}: main-new2, main-new3, null, next
  \item \texttt{addVariable}: main-new3, return
\end{itemize}

\subsubsection{\texttt{lst.splitAt(0)} に対する操作}
\begin{itemize}
  \item \texttt{editEdgeReference}: main-new1, main-new2, null, next
  \item \texttt{addVariable}: main-new2, return
\end{itemize}

\subsection{手動で復元したコード}

\subsubsection{\texttt{lst.splitAt(1)}}
\begin{lstlisting}[language=JavaScript]
var current = this;
var count = 0;
while (current !== null && count < 1) {
  current = current.next;
  count++;
}
if (current && current.next) {
  var newHead = current.next;
  current.next = null;
  return newHead;
}
return null;
\end{lstlisting}

\subsection{一般化と部分関数の推定}

\begin{lstlisting}[language=JavaScript]
// f() = 指定されたインデックスまでリストを走査し、エッジを操作して分割する関数
return f();
\end{lstlisting}

\subsection{想定の合成された関数}
\begin{lstlisting}
splitAt(i) {
  if (i < 0) return null; // Or throw error
  let current = this;
  let count = 0;

// Handle splitting at index 0 (return the whole list, make original head null)
  if (i === 0) {
    const fullList = this;
    // How to represent 'this' becoming null or a new empty node in Kanon?
    // Maybe the 'addVariable: null, return' for the original list is enough.
    return fullList;
  }

while (current !== null && count < i - 1) {
    current = current.next;
    count++;
  }

if (current !== null && current.next !== null) {
    const newHead = current.next;
    current.next = null; // editEdgeReference: current, newHead, null, next
    return newHead;
  }
  return null; // No split possible
}
\end{lstlisting}

\subsection{oracleの関数}
\begin{lstlisting}[language=JavaScript]
splitAt(i) {
  function find(node, index) {
    if (node === null || index < 0) {
      return null;
    }
    if (index === 0) {
      return node;
    }
    return find(node.next, index - 1);
  }

if (i === 0) {
    const originalList = this.next; // Assuming current 'this' becomes isolated/emptied
    this.next = null;
    return originalList;
  }

const nodeBeforeSplit = find(this, i - 1);
  if (nodeBeforeSplit !== null && nodeBeforeSplit.next !== null) {
    const newHead = nodeBeforeSplit.next;
    nodeBeforeSplit.next = null;
    return newHead;
  }
  return null;
}
\end{lstlisting}

% \clearpage

\section{対象関数: \texttt{sublist(startIndex, endIndex)}}
連結リストの指定された範囲のノードを含む新しい連結リストを作成し返すメソッド。

\subsection{初期プログラム}

\begin{lstlisting}[language=JavaScript]
class Node {
  sublist(startIndex, endIndex) { /* TBD */ }
}
var lst = new Node(); 
lst.val = 1;
lst.next = new Node();
lst.next.val = 2;
lst.next.next = new Node();
lst.next.next.val = 3;
lst.next.next.next = new Node();
lst.next.next.next.val = 4;
var sub1 = lst.sublist(0, 1); // 1, 2
var sub2 = lst.sublist(2, 3); // 3, 4
var sub3 = lst.sublist(0, 0); // 1
\end{lstlisting}

\subsection{操作ログ}

\subsubsection{\texttt{lst.sublist(0, 1)} に対する操作}
\begin{itemize}
  \item \texttt{addNode}: newSublistHead, Node, false, undefined
  \item \texttt{addNode}: newSublistHeadVal, 1, true, string
  \item \texttt{addEdge}: newSublistHead, newSublistHeadVal, val
  \item \texttt{addNode}: newSublistSecond, Node, false, undefined
  \item \texttt{addNode}: newSublistSecondVal, 2, true, string
  \item \texttt{addEdge}: newSublistSecond, newSublistSecondVal, val
  \item \texttt{addEdge}: newSublistHead, newSublistSecond, next
  \item \texttt{addVariable}: newSublistHead, return
\end{itemize}

\subsubsection{\texttt{lst.sublist(2, 3)} に対する操作}
\begin{itemize}
  \item \texttt{addNode}: newSublistHead2, Node, false, undefined
  \item \texttt{addNode}: newSublistHead2Val, 3, true, string
  \item \texttt{addEdge}: newSublistHead2, newSublistHead2Val, val
  \item \texttt{addNode}: newSublistSecond2, Node, false, undefined
  \item \texttt{addNode}: newSublistSecond2Val, 4, true, string
  \item \texttt{addEdge}: newSublistSecond2, newSublistSecond2Val, val
  \item \texttt{addEdge}: newSublistHead2, newSublistSecond2, next
  \item \texttt{addVariable}: newSublistHead2, return
\end{itemize}

\subsection{手動で復元したコード}

\subsubsection{\texttt{lst.sublist(0, 1)}}
\begin{lstlisting}[language=JavaScript]
var newList = new Node();
newList.val = this.val;
newList.next = new Node();
newList.next.val = this.next.val;
return newList;
\end{lstlisting}

\subsection{一般化と部分関数の推定}

\begin{lstlisting}[language=JavaScript]
// f() = 指定された範囲のノードを複製して新しいリストを作成する関数
return f();
\end{lstlisting}

\subsection{想定の合成された関数}
\begin{lstlisting}
sublist(startIndex, endIndex) {
  if (startIndex < 0 || startIndex > endIndex) return null; // Invalid range
  let newList = null;
  let currentNewNode = null;
  let currentOriginalNode = this;
  let currentIndex = 0;

// Advance to startIndex
  while (currentOriginalNode !== null && currentIndex < startIndex) {
    currentOriginalNode = currentOriginalNode.next;
    currentIndex++;
  }

// Copy nodes within the range
  while (currentOriginalNode !== null && currentIndex <= endIndex) {
    const newNode = new Node();
    newNode.val = currentOriginalNode.val; // addNode for newNode, addNode for val, addEdge for val

if (newList === null) {
      newList = newNode;
      currentNewNode = newNode;
    } else {
      currentNewNode.next = newNode; // addEdge for next
      currentNewNode = newNode;
    }
    currentOriginalNode = currentOriginalNode.next;
    currentIndex++;
  }
  return newList;
}
\end{lstlisting}

\subsection{oracleの関数}
\begin{lstlisting}[language=JavaScript]
sublist(startIndex, endIndex) {
  if (startIndex < 0 || startIndex > endIndex) {
    return null;
  }
  let newList = null;
  let tail = null;
  let current = this;
  let currentIndex = 0;

while (current !== null && currentIndex <= endIndex) {
    if (currentIndex >= startIndex) {
      const newNode = new Node();
      newNode.val = current.val;
      if (newList === null) {
        newList = newNode;
        tail = newNode;
      } else {
        tail.next = newNode;
        tail = newNode;
      }
    }
    current = current.next;
    currentIndex++;
  }
  return newList;
}
\end{lstlisting}

% \clearpage

\section{対象関数: \texttt{removeOdds()}}
連結リストから偶数番目のノードを削除するメソッド。

\subsection{初期プログラム}

\begin{lstlisting}[language=JavaScript]
class Node {
  removeOdds() { /* TBD */ }
}
var lst = new Node(); 
lst.val = 1; // Index 0 (keep)
lst.next = new Node();
lst.next.val = 2; // Index 1 (remove)
lst.next.next = new Node();
lst.next.next.val = 3; // Index 2 (keep)
lst.next.next.next = new Node();
lst.next.next.next.val = 4; // Index 3 (remove)
lst.removeOdds();
\end{lstlisting}

\subsection{操作ログ}

\subsubsection{\texttt{lst.removeOdds()} に対する操作}
\begin{itemize}
  \item \texttt{deleteNode}: main-new2
  \item \texttt{editEdgeReference}: main-new1, main-new2, main-new3, next
  \item \texttt{deleteNode}: main-new4
  \item \texttt{editEdgeReference}: main-new3, main-new4, null, next
\end{itemize}

\subsection{手動で復元したコード}

\subsubsection{\texttt{lst.removeOdds()}}
\begin{lstlisting}[language=JavaScript]
var current = this;
var index = 0;
while (current !== null && current.next !== null) {
  if (index % 2 === 1) { // Check if the next node is odd-indexed (0-based)
    // This logic removes the current node if it's odd.
    // If aiming to remove the node at odd index, adjust.
    // Assuming removing the NEXT node if it's odd-indexed
    current.next = current.next.next;
    // deleteNode for current.next
  } else {
    current = current.next;
  }
  index++;
}
\end{lstlisting}

\subsection{一般化と部分関数の推定}

\begin{lstlisting}[language=JavaScript]
// f() = リストを走査し、偶数番目のノードを削除する関数
f();
\end{lstlisting}

\subsection{想定の合成された関数}
\begin{lstlisting}
removeOdds() {
  if (this === null) return;
  let current = this;
  let index = 0;

// Handle the case where the head itself needs to be removed (if it's at an odd index, which is not the case for 0-indexed odd numbers)
  // Assuming 0-indexed: remove nodes at index 1, 3, 5...
  while (current !== null && current.next !== null) {
    if (index % 2 === 0) { // If current node is at an even index
      current.next = current.next.next; // Skip the next (odd-indexed) node (editEdgeReference, deleteNode)
      // current remains the same, as its next has changed
    } else {
      current = current.next; // Move to the next node
    }
    index++;
  }
}
\end{lstlisting}

\subsection{oracleの関数}
\begin{lstlisting}[language=JavaScript]
removeOdds() {
  if (this === null || this.next === null) {
    return;
  }
  // Recursive approach:
  // 'this' is at index 0 (even), so we keep it.
  // We then call removeOdds on this.next to handle the rest of the list.
  // If this.next was at index 1 (odd) and needs to be removed,
  // this.next will be replaced with this.next.next from the recursive call.
  this.next = this.next.removeOddsRecursive(1); // Helper to handle indices
}

// Helper function that explicitly takes the current index
removeOddsRecursive(currentIndex) {
  if (this === null) {
    return null;
  }
  if (currentIndex % 2 === 1) { // If current node is at an odd index
    // This node should be removed. Return its next node.
    // The caller will then link to this.next.
    return this.next ? this.next.removeOddsRecursive(currentIndex + 1) : null;
  } else {
    // This node should be kept. Link its next to the result of recursive call.
    this.next = this.next ? this.next.removeOddsRecursive(currentIndex + 1) : null;
    return this;
  }
}
\end{lstlisting}

% \clearpage

\section{対象関数: \texttt{pad(length, value)}}
連結リストを指定された\texttt{length}になるまで\texttt{value}でパディングするメソッド。

\subsection{初期プログラム}

\begin{lstlisting}[language=JavaScript]
class Node {
  pad(length, value) { /* TBD */ }
}
var lst = new Node(); 
lst.val = 1;
lst.next = new Node();
lst.next.val = 2;
lst.pad(5, 99); // Pad to length 5 with 99s
var shortLst = new Node();
shortLst.val = 1;
shortLst.pad(1, 0); // No padding needed
\end{lstlisting}

\subsection{操作ログ}

\subsubsection{\texttt{lst.pad(5, 99)} に対する操作}
\begin{itemize}
  \item \texttt{addNode}: temp1, Node, false, undefined
  \item \texttt{addNode}: temp2, 99, true, string
  \item \texttt{addEdge}: temp1, temp2, val
  \item \texttt{editEdgeReference}: main-new2, null, temp1, next
  \item \texttt{addNode}: temp3, Node, false, undefined
  \item \texttt{addNode}: temp4, 99, true, string
  \item \texttt{addEdge}: temp3, temp4, val
  \item \texttt{addEdge}: temp1, temp3, next
  \item \texttt{addNode}: temp5, Node, false, undefined
  \item \texttt{addNode}: temp6, 99, true, string
  \item \texttt{addEdge}: temp5, temp6, val
  \item \texttt{addEdge}: temp3, temp5, next
\end{itemize}

\subsubsection{\texttt{shortLst.pad(1, 0)} に対する操作}
\begin{itemize}
  \item (No operations)
\end{itemize}

\subsection{手動で復元したコード}

\subsubsection{\texttt{lst.pad(5, 99)}}
\begin{lstlisting}[language=JavaScript]
var current = this;
var count = 0;
while (current.next !== null) {
  current = current.next;
  count++;
}
count++; // Count current node

while (count < 5) {
  var newNode = new Node();
  newNode.val = 99;
  current.next = newNode;
  current = newNode;
  count++;
}
\end{lstlisting}

\subsection{一般化と部分関数の推定}

\begin{lstlisting}[language=JavaScript]
// f() = リストの長さを計算し、不足しているノードを追加する関数
f();
\end{lstlisting}

\subsection{想定の合成された関数}
\begin{lstlisting}
pad(length, value) {
  let current = this;
  let count = 0;

// Find the current length
  while (current !== null && current.next !== null) {
    current = current.next;
    count++;
  }
  // Add 1 for the head node itself
  count++;

// Add padding nodes
  while (count < length) {
    const newNode = new Node();
    newNode.val = value; // addNode for newNode, addNode for val, addEdge for val
    current.next = newNode; // addEdge for next
    current = newNode;
    count++;
  }
}
\end{lstlisting}

\subsection{oracleの関数}
\begin{lstlisting}[language=JavaScript]
pad(length, value) {
  let currentLength = 0;
  let current = this;
  while (current !== null) {
    currentLength++;
    if (current.next === null) {
      break; // Found the tail
    }
    current = current.next;
  }

while (currentLength < length) {
    const newNode = new Node();
    newNode.val = value;
    if (current === null) { // For an empty list
      // This scenario is tricky with 'this' as context.
      // Assuming 'this' is always a valid node even if it represents an empty list
      // and new nodes are appended to it.
      // If 'this' is the only node and its val is undefined, set its val and then append.
      // For simplicity, we assume this.next is where new nodes are added.
    } else {
      current.next = newNode;
      current = newNode;
    }
    currentLength++;
  }
}
\end{lstlisting}

% \clearpage

\section{対象関数: \texttt{merge(otherList)}}
別の連結リストを現在のリストと結合し、ソートされた新しい連結リストを返すメソッド（破壊的ではない場合）。

\subsection{初期プログラム}

\begin{lstlisting}[language=JavaScript]
class Node {
  merge(otherList) { /* TBD */ }
}
var lst1 = new Node(); 
lst1.val = 1;
lst1.next = new Node();
lst1.next.val = 3;
var lst2 = new Node();
lst2.val = 2;
lst2.next = new Node();
lst2.next.val = 4;
var mergedList = lst1.merge(lst2);
\end{lstlisting}

\subsection{操作ログ}

\subsubsection{\texttt{lst1.merge(lst2)} に対する操作}
\begin{itemize}
  \item \texttt{addNode}: mergedHead, Node, false, undefined
  \item \texttt{addNode}: mergedHeadVal, 1, true, string
  \item \texttt{addEdge}: mergedHead, mergedHeadVal, val
  \item \texttt{addNode}: mergedSecond, Node, false, undefined
  \item \texttt{addNode}: mergedSecondVal, 2, true, string
  \item \texttt{addEdge}: mergedSecond, mergedSecondVal, val
  \item \texttt{addEdge}: mergedHead, mergedSecond, next
  \item \texttt{addNode}: mergedThird, Node, false, undefined
  \item \texttt{addNode}: mergedThirdVal, 3, true, string
  \item \texttt{addEdge}: mergedThird, mergedThirdVal, val
  \item \texttt{addEdge}: mergedSecond, mergedThird, next
  \item \texttt{addNode}: mergedFourth, Node, false, undefined
  \item \texttt{addNode}: mergedFourthVal, 4, true, string
  \item \texttt{addEdge}: mergedFourth, mergedFourthVal, val
  \item \texttt{addEdge}: mergedThird, mergedFourth, next
  \item \texttt{addVariable}: mergedHead, return
\end{itemize}

\subsection{手動で復元したコード}

\subsubsection{\texttt{lst1.merge(lst2)}}
\begin{lstlisting}[language=JavaScript]
var mergedHead = new Node();
var currentMerged = mergedHead;
var current1 = this;
var current2 = otherList;

while (current1 !== null && current2 !== null) {
  var newNode = new Node();
  if (current1.val <= current2.val) {
    newNode.val = current1.val;
    current1 = current1.next;
  } else {
    newNode.val = current2.val;
    current2 = current2.next;
  }
  currentMerged.next = newNode;
  currentMerged = newNode;
}

while (current1 !== null) {
  var newNode = new Node();
  newNode.val = current1.val;
  currentMerged.next = newNode;
  currentMerged = newNode;
  current1 = current1.next;
}

while (current2 !== null) {
  var newNode = new Node();
  newNode.val = current2.val;
  currentMerged.next = newNode;
  currentMerged = newNode;
  current2 = current2.next;
}
return mergedHead.next; // Skip dummy head
\end{lstlisting}

\subsection{一般化と部分関数の推定}

\begin{lstlisting}[language=JavaScript]
// f() = 2つのソート済みリストをマージして新しいソート済みリストを構築する関数
return f();
\end{lstlisting}

\subsection{想定の合成された関数}
\begin{lstlisting}
merge(otherList) {
  let mergedHead = null;
  let currentMerged = null;
  let current1 = this;
  let current2 = otherList;

while (current1 !== null && current2 !== null) {
    let newNode = new Node();
    // Add newNode, its val, and edge to val
    if (current1.val <= current2.val) {
      newNode.val = current1.val;
      current1 = current1.next;
    } else {
      newNode.val = current2.val;
      current2 = current2.next;
    }

if (mergedHead === null) {
      mergedHead = newNode;
      currentMerged = newNode;
    } else {
      currentMerged.next = newNode; // addEdge for next
      currentMerged = newNode;
    }
  }

// Append remaining nodes
  let remaining = current1 !== null ? current1 : current2;
  while (remaining !== null) {
    let newNode = new Node();
    newNode.val = remaining.val;
    if (mergedHead === null) {
      mergedHead = newNode;
      currentMerged = newNode;
    } else {
      currentMerged.next = newNode;
      currentMerged = newNode;
    }
    remaining = remaining.next;
  }
  return mergedHead;
}
\end{lstlisting}

\subsection{oracleの関数}
\begin{lstlisting}[language=JavaScript]
merge(otherList) {
  if (this === null) return otherList;
  if (otherList === null) return this;

let resultHead;
  if (this.val <= otherList.val) {
    resultHead = this;
    resultHead.next = this.next ? this.next.merge(otherList) : otherList;
  } else {
    resultHead = otherList;
    resultHead.next = otherList.next ? this.merge(otherList.next) : this;
  }
  return resultHead;
}
\end{lstlisting}

% \clearpage

\section{対象関数: \texttt{compact()}}
連結リストからnullまたはundefinedの値を削除するメソッド。

\subsection{初期プログラム}

\begin{lstlisting}[language=JavaScript]
class Node {
  compact() { /* TBD */ }
}
var lst = new Node(); 
lst.val = 1;
lst.next = new Node();
lst.next.val = undefined; // Null value
lst.next.next = new Node();
lst.next.next.val = 2;
lst.next.next.next = new Node();
lst.next.next.next.val = null; // Another null value
lst.compact();
\end{lstlisting}

\subsection{操作ログ}

\subsubsection{\texttt{lst.compact()} に対する操作}
\begin{itemize}
  \item \texttt{deleteNode}: main-new2
  \item \texttt{editEdgeReference}: main-new1, main-new2, main-new3, next
  \item \texttt{deleteNode}: main-new4
  \item \texttt{editEdgeReference}: main-new3, main-new4, null, next
\end{itemize}

\subsection{手動で復元したコード}

\subsubsection{\texttt{lst.compact()}}
\begin{lstlisting}[language=JavaScript]
var current = this;
while (current !== null && current.next !== null) {
  if (current.next.val === null || current.next.val === undefined) {
    current.next = current.next.next; // Skip the null/undefined node
  } else {
    current = current.next;
  }
}
// Handle head if it's null/undefined (requires return value)
\end{lstlisting}

\subsection{一般化と部分関数の推定}

\begin{lstlisting}[language=JavaScript]
// f() = リストを走査し、null/undefinedノードを削除する関数
f();
\end{lstlisting}

\subsection{想定の合成された関数}
\begin{lstlisting}
compact() {
  // Handle head node if it's null/undefined
  // This typically means the function needs to return the new head
  // For simplicity, if 'this' is null/undefined, it might be implicitly handled by caller.
  // Assuming 'this' is always a valid starting node.
  if (this === null) return null;

let current = this;
  while (current !== null && current.next !== null) {
    if (current.next.val === null || current.next.val === undefined) {
      // Delete current.next (editEdgeReference to skip, deleteNode)
      current.next = current.next.next;
    } else {
      current = current.next;
    }
  }
  // If the original head was null/undefined, it also needs to be handled
  // For in-place modification of 'this', this is tricky.
  // If the intention is to modify the list pointed to by 'this'
  // and return the potentially new head.
  if (this.val === null || this.val === undefined) {
    return this.next ? this.next.compact() : null; // Recursively compact from next node
  }
  return this;
}
\end{lstlisting}

\subsection{oracleの関数}
\begin{lstlisting}[language=JavaScript]
compact() {
  if (this === null) {
    return null;
  }
  // If current node's value is null/undefined, return its next (effectively removing it)
  if (this.val === null || this.val === undefined) {
    return this.next ? this.next.compact() : null;
  }
  // Otherwise, keep current node and compact the rest of the list
  this.next = this.next ? this.next.compact() : null;
  return this;
}
\end{lstlisting}


\end{document}
